% SGA 3, Expose V, Loop 1 English reconstruction.
% Scope: Theorem 4.1 proof parts c), d), and opening of e), including
% Lemma 4.1.2 and its proof; stops before the proof of Theorem 4.1(iii).
% Authority: Polo--Gille Exp5-13oct24.pdf, component-local pp. 13--14
% (combined-reader pp. 275--276). Re-edited page headers 263--264;
% internal original-page markers advance through 266.
% Jacob Reinhold's CC BY 4.0 Markdown is comparison/drafting material only.

\medskip
\phantomsection\label{proof:sga3-V-4.1-c}%
\noindent\textbf{c) \hyperref[thm:sga3-V-4.1-i]{(i)} 的证明.}
令 \(Y=\operatorname{Spec}B\), \(p=\operatorname{Spec}i\),
其中 \(i:B\hookrightarrow A_0\) 为包含映射.
由 \hyperref[proof:sga3-V-4.1-a]{a)}, 态射 \(p:X_0\to Y\) 满射.
先证明 \((Y,p)\) 是所有赋环空间构成的范畴中 \((d_0,d_1)\) 的余核.
事实上, \hyperref[proof:sga3-V-4.1-b]{b)} 表明,
\(\operatorname{Spec}B\) 的底集, 是在 \(X_0\) 的底集中
将满足存在 \(z\in X_1\) 使
\[
  d_1z=y,\qquad d_0z=x.
\]
的点 \(x\) 与 \(y\) 等同而得到的.
此外, 由于 \(i\) 是整的, \(p=\operatorname{Spec}i\) 为闭映射,
所以 \(Y\) 带有由 \(X_0\) 诱导的商拓扑.

由此得到 \(p\) 是开映射. 事实上, 设 \(U'\) 为 \(X_0\) 的任意开子集.
由于 \(d_1\) 满射、有限且局部自由, 故它忠实平坦且为有限呈示,
因而为开映射. 于是
\[
  U=d_1\!\left(d_0^{-1}(U')\right),
\]
即 \(U'\) 关于 \(X_\ast\) 所定义的等价关系的饱和化, 是开集.
由于 \(Y\) 带有商拓扑, \(p(U')=p(U)\) 为开集.

最后, \(B\) 的选取以及 \(p,d_0\) 和 \(d_1\) 仿射这一事实,
表明环层的典范序列
\[
  \mathcal O_Y\longrightarrow p_\ast\mathcal O_{X_0}
  \underset{p_\ast(\delta_0)}{\overset{p_\ast(\delta_1)}
    {\rightrightarrows}}
  p_\ast d_{0\ast}\mathcal O_{X_1}
  =p_\ast d_{1\ast}\mathcal O_{X_1}
\]
正合.

只须证明 \((Y,p)\) 也是概形范畴中, 更一般地说,
局部赋环空间范畴中 \((d_0,d_1)\) 的余核.
设 \(q:X_0\to Z\) 是满足 \(qd_0=qd_1\) 的概形态射.
由上面的论证, 存在唯一的赋环空间态射 \(r:Y\to Z\), 使 \(q=rp\).
须证明对每个 \(y\in Y\), \(r\) 诱导的同态
\[
  \mathcal O_{Z,r(y)}\longrightarrow\mathcal O_{Y,y}
\]
为局部同态. 由于 \(p\) 满射, 存在某个 \(x\in X_0\) 使 \(y=p(x)\),
而此断言由 \(q\) 诱导的同态
\[
  \mathcal O_{Z,q(x)}\longrightarrow\mathcal O_{X_0,x}
\]
是局部同态这一事实得到.

\medskip
\phantomsection\label{proof:sga3-V-4.1-d}%
\noindent\textbf{d) \hyperref[thm:sga3-V-4.1-ii]{(ii)} 的证明.}
这由 \hyperref[proof:sga3-V-4.1-a]{a)} 和~\hyperref[proof:sga3-V-4.1-c]{c)} 得到.

\medskip
\phantomsection\label{proof:sga3-V-4.1-e}%
\noindent\textbf{e) \hyperref[thm:sga3-V-4.1-iii]{(iii)} 的证明.}
回顾一下, \(\underline P\) 表示概形 \(P\) 的底集,
\(\underline d:\underline P\to\underline Q\) 表示态射 \(d:P\to Q\)
诱导的映射.

\medskip
\phantomsection\label{lem:sga3-V-4.1.2}%
\noindent\textbf{引理 4.1.2.}---%
\footnote{\label{note:sga3-V-22}编者注:
这里补入了 \hyperref[lem:sga3-V-4.1.2]{本引理};
它在本篇和后续各篇 (\(\mathrm{VI_A}\), \(\mathrm{VI_B}\)) 中反复使用.
它在 SGA~3 原版 1965 年版中是引理 \(\mathrm{VI_B}\), 4.5.1.}
\textit{设 \((A,\mathfrak m)\) 为局部环, \(k\) 为其剩余域,
\(K/k\) 为域扩张. 则存在一个局部且平坦的 \(A\)-代数 \(B\),
使 \(B/\mathfrak mB\) 与 \(K\) 是 \(k\)-同构的.
若 \(K/k\) 有限, 还可选取 \(B\) 使其在 \(A\) 上有限且自由.}

这已在 EGA \(0_{\mathrm{III}}\), 10.3.1 中证明,
那里还证明了当 \(A\) 为 Noether 环时, 可选取 \(B\) 也为 Noether 环.
为方便读者, 这里回顾其证明.

令 \(A'=A[T]\), 其中 \(T\) 为不定元.
若 \(K=k(T)\), 令 \(\mathfrak p=\mathfrak mA'\), \(B=A'_{\mathfrak p}\).
则
\[
  B/\mathfrak mB\simeq k[T]_{(0)}=k(T).
\]
环 \(B\) 在 \(A'\) 上平坦, 而后者为自由 \(A\)-模,
所以 \(B\) 在 \(A\) 上平坦.

再假设 \(K=k(t)=k[t]\), 其中 \(t\) 在 \(k\) 上代数.
令 \(B=A'/(F)\), 其中 \(F\in A'\) 为首一多项式,
它在 \(k[T]\) 中的像为 \(t\) 在 \(k\) 上的极小多项式 \(f\).
则 \(B\) 是有限秩自由 \(A\)-模, 其秩为
\[
  \deg(F)=\deg(f).
\]
特别地, \(B\) 在 \(A\) 上整,
因而 \(B\) 的每个极大理想都包含 \(\mathfrak m\).
由于
\[
  B/\mathfrak mB\simeq k[T]/(f)\simeq K,
\]
\(B\) 是局部环, 其极大理想为 \(\mathfrak mB\).
这已证明当 \([K:k]<\infty\) 时, \(B\) 可以选为在 \(A\) 上有限且自由.

在一般情形中, 选取 \(K/k\) 的一族生成元 \((t_i)_{i\in I}\),
并在 \(I\) 上赋予良序; 即在 \(I\) 上赋予全序 \(\leq\),
使每个非空子集都有最小元素.
对 \(i\in I\), 令 \(k_i\) (分别地, \(k_{<i}\)) 为
\(K\) 中由满足 \(j\leq i\) (分别地, \(j<i\)) 的 \(t_j\)
生成的子域. 必要时添加一个元素, 可假设 \(I\) 有最大元素 \(\xi\),
从而 \(K=k_\xi\).

设 \(J\subseteq I\) 为使下述条件成立的指标 \(i\) 组成的集合:
存在局部且平坦的 \(A\)-代数的正向系 \((A_j)_{j\leq i}\), 满足
\[
  A_j/\mathfrak mA_j\simeq k_j
\]
且当 \(\ell<j\) 时, \(A_j\) 在 \(A_\ell\) 上平坦.
假设 \(I\setminus J\) 非空, 取其最小元素 \(i\), 并令
\[
  A'=\varinjlim_{j<i}A_j.
\]
由于张量积与正向极限可交换, \(A'\) 在 \(A\) 及每个
\(A_j\), \(j<i\), 上平坦, 并且
\[
  A'/\mathfrak mA'
  \simeq A'\otimes_A(A/\mathfrak m)
  \simeq k_{<i}.
\]
此外, \(A'\) 是局部环, 其极大理想为 \(\mathfrak mA'\).
事实上, 若 \(x=f_j(x_j)\) 不可逆,
则 \(x_j\) 不可逆, 因而属于 \(A_j\) 的极大理想 \(\mathfrak mA_j\);
所以 \(x\in\mathfrak mA'\).

由上面处理的单生成情形, 得到一个局部且平坦的 \(A'\)-代数 \(A_i\), 满足
\[
  A_i/\mathfrak mA_i\simeq k_{<i}(t_i)=k_i.
\]
代数 \(A_i\) 在每个 \(A_j\), \(j<i\), 上平坦,
所以 \(i\in J\), 与 \(i\) 的选择矛盾.
这一矛盾证明 \(J=I\), 而 \(A_\xi\) 具有所要求的性质.
这就证明了引理~\hyperref[lem:sga3-V-4.1.2]{4.1.2}.
% Translation statistics: source_lines=157; output_lines=153; inline_math=133; display_math=12; equation=0; equation_star=0; align=0; footnotes=1.
% Review: complete assigned component; pending_TODO=0; all math symbols, labels, references, diagram paths and environments checked; formula prose translated.
